% --- Archon physics formalization source begin ---
% archon:physics
% archon:covers PhyXMiniProblems/problem_phyx_mini_0796.lean
% archon:source-report reports/phyx_mini/problem_phyx_mini_0796.source.json

\chapter{Physics problem phyx\_mini\_0796}
\label{ch:PhyXMiniProblems_problem_phyx_mini_0796}

\paragraph{Problem source.}
\#\# Physical scenario

a projectile launched with speed \textbackslash{}( v\_0 \textbackslash{}) at an initial angle \textbackslash{}( \textbackslash{}alpha\_0 \textbackslash{}) between 0 and \textbackslash{}( 90\textasciicircum{}\{\textbackslash{}circ\} \textbackslash{}) as shown in figure.

\#\# Question

For a given \textbackslash{}( v\_0 \textbackslash{}), what value gives maximum horizontal range?

\#\# Answer choices

- A: A: \textbackslash{}( 30\textasciicircum{}\{\textbackslash{}circ\} \textbackslash{})

- B: B: \textbackslash{}( 45\textasciicircum{}\{\textbackslash{}circ\} \textbackslash{})

- C: C: \textbackslash{}( 60\textasciicircum{}\{\textbackslash{}circ\} \textbackslash{})

- D: D: \textbackslash{}( 90\textasciicircum{}\{\textbackslash{}circ\} \textbackslash{})

\#\# Figure caption (auxiliary; use the image as primary evidence)

The image depicts a projectile motion scenario on a two-dimensional graph with horizontal (x) and vertical (y) axes.

1. **Initial Conditions:**
   - The projectile is launched from the origin (0,0).
   - Initial velocity \textbackslash{}( v\_0 = 37.0 \textbackslash{}, \textbackslash{}text\{m/s\} \textbackslash{}).
   - Launch angle \textbackslash{}( \textbackslash{}alpha\_0 = 53.1\textasciicircum{}\textbackslash{}circ \textbackslash{}) from the horizontal.

2. **Path and Points:**
   - The projectile follows a parabolic path indicated by a dashed line.
   - At the initial position, the components \textbackslash{}( x \textbackslash{}) and \textbackslash{}( y \textbackslash{}) are unknowns (\textbackslash{}( x=? \textbackslash{}), \textbackslash{}( y=? \textbackslash{})).
   - A marked point on the curve at \textbackslash{}( t=2.00 \textbackslash{}, \textbackslash{}text\{s\} \textbackslash{}).

3. **Mid-trajectory:**
   - There is a highest point (apex) of the projectile, marked with unknowns for time \textbackslash{}( t\_1=? \textbackslash{}), velocity \textbackslash{}( v\_1 \textbackslash{}), and maximum height \textbackslash{}( h=? \textbackslash{}).

4. **Final Conditions:**
   - The projectile lands at another point on the x-axis, where the range \textbackslash{}( R=? \textbackslash{}) and time of flight \textbackslash{}( t\_2=? \textbackslash{}) are marked as unknowns.
   - Final velocity at landing is indicated as \textbackslash{}( v\_2 \textbackslash{}).

The relationships involve determining trajectory details including height, time of flight, and range based on initial conditions and the motion equations.

\#\# Dataset metadata

Kinematics; Physical Model Grounding Reasoning, Multi-Formula Reasoning

\paragraph{Recorded answer/context.}
B: B: \textbackslash{}( 45\textasciicircum{}\{\textbackslash{}circ\} \textbackslash{})

\paragraph{Figure/image path.}
/root/proposal\_for\_physic/science-mango/phyx\_mini\_run/phyx\_data/test\_image/796.png

\paragraph{Formalization target.}
create a compiling Lean file with sorry bodies at `PhyXMiniProblems/problem\_phyx\_mini\_0796.lean`. The Lean declarations must preserve the physical quantities, dimensions or dimensional roles, figure labels, governing-law hypotheses, and final relation expressed by this problem.
Use Mathlib/Physlib names found through LeanExplore where available. If a physics API is missing, introduce faithful local abstractions rather than scalar placeholder aliases.

\begin{theorem}[Physics formalization target]
\label{thm:physics:phyx_mini_0796:target}
\lean{PhyXMiniProblems.ProblemPhyXMini0796.problem_phyx_mini_0796}
\uses{def:physics:phyx-mini-0796:phyxminiproblems-problemphyxmini0796-degrees, def:physics:phyx-mini-0796:phyxminiproblems-problemphyxmini0796-projectilerangesetup, def:physics:phyx-mini-0796:phyxminiproblems-problemphyxmini0796-matchesproblemreadouts, def:physics:phyx-mini-0796:phyxminiproblems-problemphyxmini0796-matchesprimaryfigure, def:physics:phyx-mini-0796:phyxminiproblems-problemphyxmini0796-hasphysicalprojectilebranch, def:physics:phyx-mini-0796:phyxminiproblems-problemphyxmini0796-satisfiesuniformgravityprojectilelaws, def:physics:phyx-mini-0796:phyxminiproblems-problemphyxmini0796-maximizeshorizontalrange, def:physics:phyx-mini-0796:phyxminiproblems-problemphyxmini0796-answerchoice}
The assigned autoformalize agent should translate this physics problem into Lean declarations in the covered file, with theorem and lemma proof bodies written as `by sorry`.
\end{theorem}
\begin{proof}
This is an autoformalization task, not a proof task. Produce statements that can later be proved by the physics prover without weakening the physical model.
\end{proof}

% --- Archon declaration topology begin ---
\section{Declaration topology}
\begin{definition}[velocity Dimension]
  \label{def:physics:phyx-mini-0796:phyxminiproblems-problemphyxmini0796-velocitydimension}
  \lean{PhyXMiniProblems.ProblemPhyXMini0796.velocityDimension}
  The physical dimension L T⁻¹ of velocity
\end{definition}
\begin{proof}
  This is the declared model definition of velocity Dimension.
\end{proof}
\begin{definition}[acceleration Dimension]
  \label{def:physics:phyx-mini-0796:phyxminiproblems-problemphyxmini0796-accelerationdimension}
  \lean{PhyXMiniProblems.ProblemPhyXMini0796.accelerationDimension}
  The physical dimension L T⁻² of acceleration
\end{definition}
\begin{proof}
  This is the declared model definition of acceleration Dimension.
\end{proof}
\begin{definition}[Length Quantity]
  \label{def:physics:phyx-mini-0796:phyxminiproblems-problemphyxmini0796-lengthquantity}
  \lean{PhyXMiniProblems.ProblemPhyXMini0796.LengthQuantity}
  A nonnegative physical length such as the range or maximum height
\end{definition}
\begin{proof}
  This is the declared model definition of Length Quantity.
\end{proof}
\begin{definition}[Time Quantity]
  \label{def:physics:phyx-mini-0796:phyxminiproblems-problemphyxmini0796-timequantity}
  \lean{PhyXMiniProblems.ProblemPhyXMini0796.TimeQuantity}
  A nonnegative elapsed time
\end{definition}
\begin{proof}
  This is the declared model definition of Time Quantity.
\end{proof}
\begin{definition}[Speed Quantity]
  \label{def:physics:phyx-mini-0796:phyxminiproblems-problemphyxmini0796-speedquantity}
  \lean{PhyXMiniProblems.ProblemPhyXMini0796.SpeedQuantity}
  A nonnegative launch-speed magnitude
\end{definition}
\begin{proof}
  This is the declared model definition of Speed Quantity.
\end{proof}
\begin{definition}[Acceleration Magnitude Quantity]
  \label{def:physics:phyx-mini-0796:phyxminiproblems-problemphyxmini0796-accelerationmagnitudequantity}
  \lean{PhyXMiniProblems.ProblemPhyXMini0796.AccelerationMagnitudeQuantity}
  \uses{def:physics:phyx-mini-0796:phyxminiproblems-problemphyxmini0796-accelerationdimension}
  A nonnegative magnitude of gravitational acceleration
\end{definition}
\begin{proof}
  This is the declared model definition of Acceleration Magnitude Quantity.
\end{proof}
\begin{definition}[Planar Position Quantity]
  \label{def:physics:phyx-mini-0796:phyxminiproblems-problemphyxmini0796-planarpositionquantity}
  \lean{PhyXMiniProblems.ProblemPhyXMini0796.PlanarPositionQuantity}
  A signed planar position vector
\end{definition}
\begin{proof}
  This is the declared model definition of Planar Position Quantity.
\end{proof}
\begin{definition}[Planar Velocity Quantity]
  \label{def:physics:phyx-mini-0796:phyxminiproblems-problemphyxmini0796-planarvelocityquantity}
  \lean{PhyXMiniProblems.ProblemPhyXMini0796.PlanarVelocityQuantity}
  \uses{def:physics:phyx-mini-0796:phyxminiproblems-problemphyxmini0796-velocitydimension}
  A signed planar velocity vector
\end{definition}
\begin{proof}
  This is the declared model definition of Planar Velocity Quantity.
\end{proof}
\begin{definition}[x Axis]
  \label{def:physics:phyx-mini-0796:phyxminiproblems-problemphyxmini0796-xaxis}
  \lean{PhyXMiniProblems.ProblemPhyXMini0796.xAxis}
  Coordinate 0, horizontal and positive to the right
\end{definition}
\begin{proof}
  This is the declared model definition of x Axis.
\end{proof}
\begin{definition}[y Axis]
  \label{def:physics:phyx-mini-0796:phyxminiproblems-problemphyxmini0796-yaxis}
  \lean{PhyXMiniProblems.ProblemPhyXMini0796.yAxis}
  Coordinate 1, vertical and positive upward
\end{definition}
\begin{proof}
  This is the declared model definition of y Axis.
\end{proof}
\begin{definition}[length In Meters]
  \label{def:physics:phyx-mini-0796:phyxminiproblems-problemphyxmini0796-lengthinmeters}
  \lean{PhyXMiniProblems.ProblemPhyXMini0796.lengthInMeters}
  \uses{def:physics:phyx-mini-0796:phyxminiproblems-problemphyxmini0796-lengthquantity}
  Read a physical length in metres
\end{definition}
\begin{proof}
  This is the declared model definition of length In Meters.
\end{proof}
\begin{definition}[time In Seconds]
  \label{def:physics:phyx-mini-0796:phyxminiproblems-problemphyxmini0796-timeinseconds}
  \lean{PhyXMiniProblems.ProblemPhyXMini0796.timeInSeconds}
  \uses{def:physics:phyx-mini-0796:phyxminiproblems-problemphyxmini0796-timequantity}
  Read a physical elapsed time in seconds
\end{definition}
\begin{proof}
  This is the declared model definition of time In Seconds.
\end{proof}
\begin{definition}[speed In Meters Per Second]
  \label{def:physics:phyx-mini-0796:phyxminiproblems-problemphyxmini0796-speedinmeterspersecond}
  \lean{PhyXMiniProblems.ProblemPhyXMini0796.speedInMetersPerSecond}
  \uses{def:physics:phyx-mini-0796:phyxminiproblems-problemphyxmini0796-speedquantity}
  Read a physical speed in metres per second
\end{definition}
\begin{proof}
  This is the declared model definition of speed In Meters Per Second.
\end{proof}
\begin{definition}[acceleration In Meters Per Second Squared]
  \label{def:physics:phyx-mini-0796:phyxminiproblems-problemphyxmini0796-accelerationinmeterspersecondsquared}
  \lean{PhyXMiniProblems.ProblemPhyXMini0796.accelerationInMetersPerSecondSquared}
  \uses{def:physics:phyx-mini-0796:phyxminiproblems-problemphyxmini0796-accelerationmagnitudequantity}
  Read a gravitational-acceleration magnitude in metres per second squared
\end{definition}
\begin{proof}
  This is the declared model definition of acceleration In Meters Per Second Squared.
\end{proof}
\begin{definition}[position In Meters]
  \label{def:physics:phyx-mini-0796:phyxminiproblems-problemphyxmini0796-positioninmeters}
  \lean{PhyXMiniProblems.ProblemPhyXMini0796.positionInMeters}
  \uses{def:physics:phyx-mini-0796:phyxminiproblems-problemphyxmini0796-planarpositionquantity}
  Read a planar position in Cartesian metres
\end{definition}
\begin{proof}
  This is the declared model definition of position In Meters.
\end{proof}
\begin{definition}[velocity In Meters Per Second]
  \label{def:physics:phyx-mini-0796:phyxminiproblems-problemphyxmini0796-velocityinmeterspersecond}
  \lean{PhyXMiniProblems.ProblemPhyXMini0796.velocityInMetersPerSecond}
  \uses{def:physics:phyx-mini-0796:phyxminiproblems-problemphyxmini0796-planarvelocityquantity}
  Read a planar velocity in Cartesian metres per second
\end{definition}
\begin{proof}
  This is the declared model definition of velocity In Meters Per Second.
\end{proof}
\begin{definition}[degrees]
  \label{def:physics:phyx-mini-0796:phyxminiproblems-problemphyxmini0796-degrees}
  \lean{PhyXMiniProblems.ProblemPhyXMini0796.degrees}
  Convert a degree readout to Mathlib's angle modulo 2 * π
\end{definition}
\begin{proof}
  This is the declared model definition of degrees.
\end{proof}
\begin{definition}[Projectile Model]
  \label{def:physics:phyx-mini-0796:phyxminiproblems-problemphyxmini0796-projectilemodel}
  \lean{PhyXMiniProblems.ProblemPhyXMini0796.ProjectileModel}
  The idealized mechanics model used by the textbook range question
\end{definition}
\begin{proof}
  This is the declared model definition of Projectile Model.
\end{proof}
\begin{definition}[Figure Point]
  \label{def:physics:phyx-mini-0796:phyxminiproblems-problemphyxmini0796-figurepoint}
  \lean{PhyXMiniProblems.ProblemPhyXMini0796.FigurePoint}
  Distinguished points visible along the dashed trajectory
\end{definition}
\begin{proof}
  This is the declared model definition of Figure Point.
\end{proof}
\begin{definition}[Figure Velocity Label]
  \label{def:physics:phyx-mini-0796:phyxminiproblems-problemphyxmini0796-figurevelocitylabel}
  \lean{PhyXMiniProblems.ProblemPhyXMini0796.FigureVelocityLabel}
  Velocity arrows and their literal names in the supplied figure
\end{definition}
\begin{proof}
  This is the declared model definition of Figure Velocity Label.
\end{proof}
\begin{definition}[Figure Element]
  \label{def:physics:phyx-mini-0796:phyxminiproblems-problemphyxmini0796-figureelement}
  \lean{PhyXMiniProblems.ProblemPhyXMini0796.FigureElement}
  Visible graphical elements in image 796.png
\end{definition}
\begin{proof}
  This is the declared model definition of Figure Element.
\end{proof}
\begin{definition}[Figure Label]
  \label{def:physics:phyx-mini-0796:phyxminiproblems-problemphyxmini0796-figurelabel}
  \lean{PhyXMiniProblems.ProblemPhyXMini0796.FigureLabel}
  Literal numerical and symbolic labels printed in image 796.png
\end{definition}
\begin{proof}
  This is the declared model definition of Figure Label.
\end{proof}
\begin{definition}[Projectile Figure]
  \label{def:physics:phyx-mini-0796:phyxminiproblems-problemphyxmini0796-projectilefigure}
  \lean{PhyXMiniProblems.ProblemPhyXMini0796.ProjectileFigure}
  \uses{def:physics:phyx-mini-0796:phyxminiproblems-problemphyxmini0796-lengthquantity, def:physics:phyx-mini-0796:phyxminiproblems-problemphyxmini0796-timequantity, def:physics:phyx-mini-0796:phyxminiproblems-problemphyxmini0796-planarvelocityquantity, def:physics:phyx-mini-0796:phyxminiproblems-problemphyxmini0796-figurepoint, def:physics:phyx-mini-0796:phyxminiproblems-problemphyxmini0796-figurevelocitylabel, def:physics:phyx-mini-0796:phyxminiproblems-problemphyxmini0796-figureelement, def:physics:phyx-mini-0796:phyxminiproblems-problemphyxmini0796-figurelabel}
  Qualitative and schematic information transcribed from the raster. The schematic coordinates describe the drawing only; physical coordinates, times, lengths, and velocities remain unit-aware quantities
\end{definition}
\begin{proof}
  This is the declared model definition of Projectile Figure.
\end{proof}
\begin{definition}[Projectile Range Setup]
  \label{def:physics:phyx-mini-0796:phyxminiproblems-problemphyxmini0796-projectilerangesetup}
  \lean{PhyXMiniProblems.ProblemPhyXMini0796.ProjectileRangeSetup}
  \uses{def:physics:phyx-mini-0796:phyxminiproblems-problemphyxmini0796-lengthquantity, def:physics:phyx-mini-0796:phyxminiproblems-problemphyxmini0796-timequantity, def:physics:phyx-mini-0796:phyxminiproblems-problemphyxmini0796-speedquantity, def:physics:phyx-mini-0796:phyxminiproblems-problemphyxmini0796-accelerationmagnitudequantity, def:physics:phyx-mini-0796:phyxminiproblems-problemphyxmini0796-planarpositionquantity, def:physics:phyx-mini-0796:phyxminiproblems-problemphyxmini0796-planarvelocityquantity, def:physics:phyx-mini-0796:phyxminiproblems-problemphyxmini0796-projectilemodel, def:physics:phyx-mini-0796:phyxminiproblems-problemphyxmini0796-projectilefigure}
  A fixed-speed family of same-height projectile launches. The trajectory, flight time, and horizontal range are indexed by the candidate physical launch angle. The angle shown in the supplied figure is a separate field and does not define the maximizing angle
\end{definition}
\begin{proof}
  This is the declared model definition of Projectile Range Setup.
\end{proof}
\begin{definition}[Is Admissible Launch Angle]
  \label{def:physics:phyx-mini-0796:phyxminiproblems-problemphyxmini0796-isadmissiblelaunchangle}
  \lean{PhyXMiniProblems.ProblemPhyXMini0796.IsAdmissibleLaunchAngle}
  Candidate launch angles run from the horizontal through the vertical
\end{definition}
\begin{proof}
  This is the declared model definition of Is Admissible Launch Angle.
\end{proof}
\begin{definition}[Matches Problem Readouts]
  \label{def:physics:phyx-mini-0796:phyxminiproblems-problemphyxmini0796-matchesproblemreadouts}
  \lean{PhyXMiniProblems.ProblemPhyXMini0796.MatchesProblemReadouts}
  \uses{def:physics:phyx-mini-0796:phyxminiproblems-problemphyxmini0796-xaxis, def:physics:phyx-mini-0796:phyxminiproblems-problemphyxmini0796-yaxis, def:physics:phyx-mini-0796:phyxminiproblems-problemphyxmini0796-lengthinmeters, def:physics:phyx-mini-0796:phyxminiproblems-problemphyxmini0796-timeinseconds, def:physics:phyx-mini-0796:phyxminiproblems-problemphyxmini0796-speedinmeterspersecond, def:physics:phyx-mini-0796:phyxminiproblems-problemphyxmini0796-positioninmeters, def:physics:phyx-mini-0796:phyxminiproblems-problemphyxmini0796-velocityinmeterspersecond, def:physics:phyx-mini-0796:phyxminiproblems-problemphyxmini0796-degrees, def:physics:phyx-mini-0796:phyxminiproblems-problemphyxmini0796-projectilerangesetup}
  The numerical values and physical point identifications printed in the figure. The 53.1 degree figure angle is illustrative and is not asserted to maximize the range
\end{definition}
\begin{proof}
  This is the declared model definition of Matches Problem Readouts.
\end{proof}
\begin{definition}[Matches Primary Figure]
  \label{def:physics:phyx-mini-0796:phyxminiproblems-problemphyxmini0796-matchesprimaryfigure}
  \lean{PhyXMiniProblems.ProblemPhyXMini0796.MatchesPrimaryFigure}
  \uses{def:physics:phyx-mini-0796:phyxminiproblems-problemphyxmini0796-projectilerangesetup}
  Primary-image evidence for the axes, dashed path, labelled points, and four velocity arrows. No schematic coordinate is treated as a calibrated length
\end{definition}
\begin{proof}
  This is the declared model definition of Matches Primary Figure.
\end{proof}
\begin{definition}[Has Physical Projectile Branch]
  \label{def:physics:phyx-mini-0796:phyxminiproblems-problemphyxmini0796-hasphysicalprojectilebranch}
  \lean{PhyXMiniProblems.ProblemPhyXMini0796.HasPhysicalProjectileBranch}
  \uses{def:physics:phyx-mini-0796:phyxminiproblems-problemphyxmini0796-timeinseconds, def:physics:phyx-mini-0796:phyxminiproblems-problemphyxmini0796-speedinmeterspersecond, def:physics:phyx-mini-0796:phyxminiproblems-problemphyxmini0796-accelerationinmeterspersecondsquared, def:physics:phyx-mini-0796:phyxminiproblems-problemphyxmini0796-degrees, def:physics:phyx-mini-0796:phyxminiproblems-problemphyxmini0796-projectilerangesetup, def:physics:phyx-mini-0796:phyxminiproblems-problemphyxmini0796-isadmissiblelaunchangle}
  Positivity and endpoint-branch conditions. At a horizontal launch the same-height flight is the zero-duration endpoint; every admissible angle strictly above the horizontal uses the later, positive landing time
\end{definition}
\begin{proof}
  This is the declared model definition of Has Physical Projectile Branch.
\end{proof}
\begin{definition}[Satisfies Uniform Gravity Projectile Laws]
  \label{def:physics:phyx-mini-0796:phyxminiproblems-problemphyxmini0796-satisfiesuniformgravityprojectilelaws}
  \lean{PhyXMiniProblems.ProblemPhyXMini0796.SatisfiesUniformGravityProjectileLaws}
  \uses{def:physics:phyx-mini-0796:phyxminiproblems-problemphyxmini0796-xaxis, def:physics:phyx-mini-0796:phyxminiproblems-problemphyxmini0796-yaxis, def:physics:phyx-mini-0796:phyxminiproblems-problemphyxmini0796-lengthinmeters, def:physics:phyx-mini-0796:phyxminiproblems-problemphyxmini0796-timeinseconds, def:physics:phyx-mini-0796:phyxminiproblems-problemphyxmini0796-speedinmeterspersecond, def:physics:phyx-mini-0796:phyxminiproblems-problemphyxmini0796-accelerationinmeterspersecondsquared, def:physics:phyx-mini-0796:phyxminiproblems-problemphyxmini0796-positioninmeters, def:physics:phyx-mini-0796:phyxminiproblems-problemphyxmini0796-velocityinmeterspersecond, def:physics:phyx-mini-0796:phyxminiproblems-problemphyxmini0796-projectilerangesetup, def:physics:phyx-mini-0796:phyxminiproblems-problemphyxmini0796-isadmissiblelaunchangle}
  The standard no-drag, constant-gravity projectile equations in coherent SI readouts. They are stated for every admissible candidate angle at the fixed launch speed. The final two fields say only that a flight ends at its launch height and that R measures its horizontal displacement; neither identifies the maximizing angle
\end{definition}
\begin{proof}
  This is the declared model definition of Satisfies Uniform Gravity Projectile Laws.
\end{proof}
\begin{lemma}[horizontal Range Formula]
  \label{lem:physics:phyx-mini-0796:phyxminiproblems-problemphyxmini0796-horizontalrangeformula}
  \lean{PhyXMiniProblems.ProblemPhyXMini0796.horizontalRangeFormula}
  \uses{def:physics:phyx-mini-0796:phyxminiproblems-problemphyxmini0796-lengthinmeters, def:physics:phyx-mini-0796:phyxminiproblems-problemphyxmini0796-speedinmeterspersecond, def:physics:phyx-mini-0796:phyxminiproblems-problemphyxmini0796-accelerationinmeterspersecondsquared, def:physics:phyx-mini-0796:phyxminiproblems-problemphyxmini0796-projectilerangesetup, def:physics:phyx-mini-0796:phyxminiproblems-problemphyxmini0796-isadmissiblelaunchangle, def:physics:phyx-mini-0796:phyxminiproblems-problemphyxmini0796-hasphysicalprojectilebranch, def:physics:phyx-mini-0796:phyxminiproblems-problemphyxmini0796-satisfiesuniformgravityprojectilelaws}
  Eliminating the positive flight time from the same-height endpoint equation gives the usual range relation. This is a derived conclusion from the model, not an assumption and not the requested maximization result
\end{lemma}
\begin{proof}
  The conclusion follows from the governing relations and figure readouts named in the statement of horizontal Range Formula.
\end{proof}
\begin{definition}[Maximizes Horizontal Range]
  \label{def:physics:phyx-mini-0796:phyxminiproblems-problemphyxmini0796-maximizeshorizontalrange}
  \lean{PhyXMiniProblems.ProblemPhyXMini0796.MaximizesHorizontalRange}
  \uses{def:physics:phyx-mini-0796:phyxminiproblems-problemphyxmini0796-lengthinmeters, def:physics:phyx-mini-0796:phyxminiproblems-problemphyxmini0796-projectilerangesetup, def:physics:phyx-mini-0796:phyxminiproblems-problemphyxmini0796-isadmissiblelaunchangle}
  A candidate angle gives at least as much range as every admissible angle
\end{definition}
\begin{proof}
  This is the declared model definition of Maximizes Horizontal Range.
\end{proof}
\begin{definition}[Answer Choice]
  \label{def:physics:phyx-mini-0796:phyxminiproblems-problemphyxmini0796-answerchoice}
  \lean{PhyXMiniProblems.ProblemPhyXMini0796.AnswerChoice}
  The four answer labels printed with the problem
\end{definition}
\begin{proof}
  This is the declared model definition of Answer Choice.
\end{proof}
\begin{definition}[degree Value]
  \label{def:physics:phyx-mini-0796:phyxminiproblems-problemphyxmini0796-answerchoice-degreevalue}
  \lean{PhyXMiniProblems.ProblemPhyXMini0796.AnswerChoice.degreeValue}
  \uses{def:physics:phyx-mini-0796:phyxminiproblems-problemphyxmini0796-answerchoice}
  Numerical degree values printed beside the answer choices
\end{definition}
\begin{proof}
  This is the declared model definition of degree Value.
\end{proof}
\begin{definition}[angle]
  \label{def:physics:phyx-mini-0796:phyxminiproblems-problemphyxmini0796-answerchoice-angle}
  \lean{PhyXMiniProblems.ProblemPhyXMini0796.AnswerChoice.angle}
  \uses{def:physics:phyx-mini-0796:phyxminiproblems-problemphyxmini0796-degrees, def:physics:phyx-mini-0796:phyxminiproblems-problemphyxmini0796-answerchoice}
  The physical launch angle represented by an answer choice
\end{definition}
\begin{proof}
  This is the declared model definition of angle.
\end{proof}
\begin{definition}[Matches Maximum Range]
  \label{def:physics:phyx-mini-0796:phyxminiproblems-problemphyxmini0796-answerchoice-matchesmaximumrange}
  \lean{PhyXMiniProblems.ProblemPhyXMini0796.AnswerChoice.MatchesMaximumRange}
  \uses{def:physics:phyx-mini-0796:phyxminiproblems-problemphyxmini0796-projectilerangesetup, def:physics:phyx-mini-0796:phyxminiproblems-problemphyxmini0796-maximizeshorizontalrange, def:physics:phyx-mini-0796:phyxminiproblems-problemphyxmini0796-answerchoice}
  An answer choice matches exactly when its angle maximizes the range
\end{definition}
\begin{proof}
  This is the declared model definition of Matches Maximum Range.
\end{proof}
% --- Archon declaration topology end ---
% --- Archon physics formalization source end ---
